\documentclass[aspectratio=169,11pt]{beamer}
\usepackage[utf8]{inputenc}
\usepackage[T1]{fontenc}
\usepackage[spanish]{babel}
\usepackage{lmodern}
\usepackage{booktabs}
\usepackage{tabularx}
\usepackage{hyperref}
\usepackage{enumitem}

\usetheme{Madrid}
\usecolortheme{default}
\setbeamertemplate{navigation symbols}{}
\setbeamertemplate{footline}[frame number]

\title{IA agéntica para la creación de contenido académico}
\subtitle{Sesión 5 --- Presentaciones I: narrativa y storyboard}
\author{Juliho Castillo}
\institute{CADi 40\,h · Escuela de Ingeniería y Ciencias\\
Tecnológico de Monterrey}
\date{Clave sugerida: \texttt{CADI-EIC-IA-AGENTICA-01}}

\begin{document}

\begin{frame}
\titlepage
\end{frame}

\begin{frame}{Agenda (3\,h) y producto E5}
\begin{tabularx}{\textwidth}{@{}lclX@{}}
\toprule
Bloque & Tiempo & Contenido \\
\midrule
Recap & 10 min & De texto largo a experiencia de clase \\
Narrativa & 30 min & Objetivos $\rightarrow$ mensajes $\rightarrow$ arco \\
Demo storyboard & 20 min & 8--12 nodos \\
Taller & 75 min & Storyboard de unidad EIC \\
Pares & 25 min & Checklist narrativa \\
Cierre & 10 min & TI hacia deck \\
\bottomrule
\end{tabularx}
\vspace{0.5em}
\textbf{E5:} storyboard narrativo + justificación de mensajes clave.
\end{frame}

\begin{frame}{Resultados de aprendizaje (sesión)}
Al finalizar, usted podrá:
\begin{enumerate}
\item Convertir objetivos de unidad en \textbf{mensajes clave} y arco narrativo
\item Producir un \textbf{storyboard} de 8--12 nodos (gancho, desarrollo, cierre)
\item Anticipar \textbf{malentendidos} típicos del tema STEM
\item Justificar \textbf{carga cognitiva} (qué va en slide vs.\ en notas)
\end{enumerate}
\end{frame}

\begin{frame}{De capítulo/sección a experiencia de clase}
\begin{itemize}
\item El texto largo no se ``pega'' en diapositivas
\item La clase es una \textbf{experiencia}: ritmo, énfasis, pausas
\item El storyboard es el contrato del deck (sesión 6 lo materializa)
\item Una idea fuerte por slide; el detalle vive en notas de orador
\end{itemize}
\begin{block}{CEDDIE}
Verificar política de integridad / uso de IA por edición.
\end{block}
\end{frame}

\begin{frame}{Objetivos $\rightarrow$ mensajes clave}
\begin{enumerate}
\item Liste 3--5 objetivos de la unidad
\item Redúzcalos a \textbf{3 mensajes} que el estudiante debe recordar
\item Cada mensaje necesita evidencia (ejemplo, ecuación, caso)
\item Lo que no sirve a un mensaje: candidato a notas o a cortar
\end{enumerate}
\end{frame}

\begin{frame}{Arco narrativo y carga cognitiva}
\begin{center}
\texttt{Gancho $\rightarrow$ Desarrollo $\rightarrow$ Práctica $\rightarrow$ Cierre}
\end{center}
\begin{itemize}
\item Gancho: por qué importa (problema EIC real)
\item Desarrollo: 1 idea por nodo; densidad controlada
\item Práctica: ejercicio breve o pregunta de chequeo
\item Cierre: síntesis + puente
\end{itemize}
\begin{alertblock}{Carga cognitiva}
Si la slide compite con usted, sobra texto.
\end{alertblock}
\end{frame}

\begin{frame}{Notación en slides}
\begin{itemize}
\item Símbolos grandes y legibles a distancia
\item Una ecuación central por slide (o paso a paso)
\item Defina símbolos la primera vez que aparecen
\item Evite muros de derivación: use notas o pizarra
\end{itemize}
\end{frame}

\begin{frame}{Anticipar malentendidos STEM}
Para cada nodo clave pregunte:
\begin{itemize}
\item ¿Qué confusión típica anticipo?
\item ¿Qué ejemplo o contraejemplo la desactiva?
\item ¿Qué debo decir en voz alta (notas de orador)?
\end{itemize}
Anticipar malentendidos = diseño instruccional, no adorno.
\end{frame}

\begin{frame}{Demo --- storyboard 8--12 nodos}
Construcción en vivo:
\begin{enumerate}
\item 3 mensajes clave desde objetivos
\item Tabla: \# | título | mensaje | visual/ecuación | notas | tiempo
\item Recortar nodos que no sirven al arco
\item \textbf{HITL:} aprobar densidad y notación
\end{enumerate}
\end{frame}

\begin{frame}{Ejemplo de tabla storyboard}
\begin{tabularx}{\textwidth}{@{}clXX@{}}
\toprule
\# & Título & Mensaje & Visual \\
\midrule
1 & Problema & Por qué importa & Foto/caso \\
2 & Notación & Símbolos del día & Tabla corta \\
3 & Idea central & Definición operativa & Ecuación \\
\ldots & \ldots & \ldots & \ldots \\
12 & Cierre & 3 takeaways & Lista \\
\bottomrule
\end{tabularx}
\vspace{0.5em}
Complete columnas de notas de orador y tiempo estimado.
\end{frame}

\begin{frame}{Qué va en slide vs.\ en notas}
\begin{tabularx}{\textwidth}{@{}XX@{}}
\toprule
\textbf{En la slide} & \textbf{En notas de orador} \\
\midrule
Mensaje + visual clave & Guion de 3--6 líneas \\
Ecuación esencial & Derivación o anécdota \\
Pregunta al grupo & Respuestas esperadas \\
\bottomrule
\end{tabularx}
\vspace{0.8em}
Sesión 6 profundiza notas y accesibilidad.
\end{frame}

\begin{frame}{Actividad de taller ($\sim$75 min)}
\textbf{Entregable E5:}
\begin{enumerate}
\item \textbf{Parte A --- Mensajes y arco} (25 min): 3 mensajes + gancho/desarrollo/cierre
\item \textbf{Parte B --- Storyboard} (35 min): tabla 8--12 filas
\item \textbf{Parte C --- Malentendidos} (10 min): confusión anticipada por nodo clave
\end{enumerate}
Nombre: \texttt{Apellido\_Sesion5\_storyboard}
\end{frame}

\begin{frame}{Rúbrica de pares}
\begin{tabularx}{\textwidth}{@{}lXX@{}}
\toprule
Criterio & Bien & A mejorar \\
\midrule
Arco & Gancho y cierre claros & Lista suelta \\
Densidad & 1 idea / nodo & Paredes de texto \\
STEM & Notación legible & Símbolos minúsculos \\
Malentendidos & Anticipados & Ignorados \\
\bottomrule
\end{tabularx}
\end{frame}

\begin{frame}{Roles agénticos en narrativa}
\begin{itemize}
\item Planificador: propone arco y nodos
\item Pedagogo: recorta mensajes y carga cognitiva
\item Humano: decide énfasis, ejemplos locales, tiempo real
\end{itemize}
No acepte un storyboard ``bonito'' que no pueda enseñar en 50 minutos.
\end{frame}

\begin{frame}{Takeaways de la sesión}
\begin{enumerate}
\item Storyboard = contrato del deck
\item Una idea fuerte por slide
\item Anticipar malentendidos es diseño
\item Carga cognitiva: slide vs.\ notas
\end{enumerate}
\end{frame}

\begin{frame}{Trabajo independiente --- bloque C}
Antes de la sesión 6:
\begin{itemize}
\item Refinar storyboard (tiempos y malentendidos)
\item Preparar herramienta de diapositivas
\item Traer E5 listo para materializar en deck MV
\end{itemize}
Parte de las \textbf{16\,h} independientes.
\end{frame}

\begin{frame}{Prep sesión 6 --- Presentaciones II}
Próxima sesión (3\,h):
\begin{itemize}
\item Deck mínimo viable (8--12 slides)
\item Notas de orador y variantes clase/conferencia
\item Mini-checklist de accesibilidad
\item Declarar microartefacto final hacia sesión 8
\end{itemize}
\end{frame}

\begin{frame}{Gracias}
\begin{center}
{\Large Juliho Castillo}\\[0.5em]
Escuela de Ingeniería y Ciencias\\[1em]
CADi · IA agéntica para la creación de contenido académico\\[1.5em]
{\small \emph{Verificar políticas de integridad y uso de IA con CEDDIE antes de cada edición.}}
\end{center}
\end{frame}

\end{document}
